% Copyright (c) 2026 Eugene Joseph M. Ragasa.
% SPDX-License-Identifier: Apache-2.0
% Author: Eugene Joseph M. Ragasa
% Formal proof completed: 2026-10-27
\documentclass[11pt]{article}
\usepackage[T1]{fontenc}\usepackage{amsmath,amssymb}\usepackage[margin=1in]{geometry}\usepackage[hidelinks]{hyperref}
\title{Proof Dossier: Joint Threshold Unresolved}\author{Eugene Joseph M. Ragasa}\date{Formal proof completed: 2026-10-27}
\begin{document}\maketitle
\section{Identity and status}
ID \texttt{BG-CERT-THM-023}; status \texttt{proved}. Lean declaration:
\par\noindent{\tiny\path{ProjectKoios.BandGauge.Certificates.}}\\
{\tiny\path{classifyThresholdFeasible_eq_unresolved_iff}}.
\section{Statement and proof}
Joint classification is unresolved exactly when no objective is rejected and at least one
objective is unresolved. Exhausting the three constructors proves both directions; the
empty-objective case remains accepted rather than unresolved.
\section{Boundary}
The theorem supplies no decisive feasibility status and no external inventory guarantee.
Kernel axioms are \texttt{propext}, \texttt{Classical.choice}, and \texttt{Quot.sound};
no admissions. Checked with Lean \texttt{4.35.0} and Mathlib revision
\texttt{56a005855ea17e0b0958f73ec801ae140b0a1679}. AI-assisted and kernel-checked;
no independent expert review recorded.
\end{document}
